\section{强化学习核心概念}

本节详细介绍强化学习中的五个核心概念：Agent（智能体）、Environment（环境）、Action（动作）、Observation/State（观测/状态）、Reward（奖励），以及它们之间的交互关系。

\subsection{Agent 与 Environment}

\begin{figure}[H]
\centering
\includegraphics[width=0.85\textwidth]{slides-images/slide-11.jpg}
\caption{RL 的核心交互回路：Agent 通过 Action 与 Environment 交互，接收 Observation\protect\footnotemark}
\end{figure}
\footnotetext{Slides 第11页。}

\begin{itemize}
    \item \textbf{Agent}（智能体）：系统中的决策主体，负责根据当前观测采取行动。例如：Super Mario 中的马里奥角色、无人机、自动驾驶汽车的控制器
    \item \textbf{Environment}（环境）：Agent 所存在并运行的世界。环境的规则定义了 Agent 的行动如何影响世界状态
\end{itemize}

\subsection{Action 与 Action Space}

\textbf{Action}（动作）是 Agent 在环境中可以执行的操作，记为 $a_t$（时刻 $t$ 的动作）。所有可能动作的集合称为 \textbf{Action Space}（动作空间）$\mathcal{A}$。

动作空间可以是：
\begin{itemize}
    \item \textbf{离散的}（Discrete）：如上下左右四个方向
    \item \textbf{连续的}（Continuous）：如方向盘的转角（任意实数值）
\end{itemize}

\subsection{Observation 与 State}

Agent 采取行动后，Environment 会返回一个新的\textbf{观测}（Observation），反映环境因该行动而产生的变化。Agent 所感知到的当前情况称为 \textbf{State}（状态），记为 $s_t$。

\begin{figure}[H]
\centering
\includegraphics[width=0.85\textwidth]{slides-images/slide-12.jpg}
\caption{状态变化：Agent 行动后，环境返回新的状态 $s_{t+1}$\protect\footnotemark}
\end{figure}
\footnotetext{Slides 第12页。}

\subsection{Reward 与 Return}

\textbf{Reward}（奖励）$r_t$ 是环境在每个时间步反馈给 Agent 的标量信号，表示当前状态的"好坏程度"。Agent 的目标是最大化\textbf{总回报}（Total Reward / Return）：

$$R_t = \sum_{i=t}^{\infty} r_i$$

\begin{itemize}
    \item $R_t$：从时刻 $t$ 开始的总回报
    \item $r_i$：时刻 $i$ 获得的即时奖励
\end{itemize}

\begin{figure}[H]
\centering
\includegraphics[width=0.85\textwidth]{slides-images/slide-14.jpg}
\caption{完整的 RL 交互回路：Agent 发送 Action，Environment 返回 State 和 Reward\protect\footnotemark}
\end{figure}
\footnotetext{Slides 第14页。}

\subsection{折扣回报（Discounted Return）}

在实际应用中，并非所有奖励都同等重要------\textbf{近期奖励}通常比远期奖励更有价值（因为未来具有不确定性）。因此引入\textbf{折扣因子} $\gamma \in [0, 1)$ 来对未来奖励进行衰减：

$$R_t = \sum_{i=t}^{\infty} \gamma^{i-t} r_i = r_t + \gamma r_{t+1} + \gamma^2 r_{t+2} + \cdots$$

\begin{itemize}
    \item $\gamma$：折扣因子（Discount Factor），控制对未来奖励的重视程度
    \item $\gamma$ 越接近 0，Agent 越"短视"，只关注即时奖励
    \item $\gamma$ 越接近 1，Agent 越"远视"，更关注长期回报
\end{itemize}

\begin{figure}[H]
\centering
\includegraphics[width=0.85\textwidth]{slides-images/slide-16.jpg}
\caption{折扣回报：用折扣因子 $\gamma$ 衰减未来奖励的权重\protect\footnotemark}
\end{figure}
\footnotetext{Slides 第16页。}

\begin{warningbox}{为什么需要折扣因子？}
如果不使用折扣因子（$\gamma = 1$），在无限时间步的情况下总回报可能趋向无穷大，导致优化问题不可解。此外，折扣因子也反映了现实世界中"一鸟在手胜过十鸟在林"的直觉------未来的奖励存在不确定性，当前的确定奖励更有价值。
\end{warningbox}

\subsection{Q-function（状态-动作价值函数）}

\textbf{Q-function} 是强化学习中最核心的概念之一。它衡量了在状态 $s$ 下执行动作 $a$ 后，Agent 所能获得的\textbf{期望未来总回报}：

$$Q(s_t, a_t) = \mathbb{E}[R_t \mid s_t, a_t]$$

\begin{figure}[H]
\centering
\includegraphics[width=0.85\textwidth]{slides-images/slide-18.jpg}
\caption{Q-function 的定义：期望未来总回报\protect\footnotemark}
\end{figure}
\footnotetext{Slides 第18页。}

\begin{importantbox}{Q-function 的核心含义}
$Q(s, a)$ 回答的问题是："如果我在状态 $s$ 下采取动作 $a$，并且此后一直采取最优策略，我能获得多少总回报？"它不仅考虑当前动作的即时奖励，还考虑了该动作对所有\textbf{未来状态和奖励}的影响。
\end{importantbox}

\subsection{Policy（策略）}

有了 Q-function，Agent 就可以制定\textbf{策略}（Policy），记为 $\pi(s)$。最优策略 $\pi^*(s)$ 是在每个状态下选择使 Q-function 最大化的动作：

$$\pi^*(s) = \arg\max_a Q(s, a)$$

\begin{figure}[H]
\centering
\includegraphics[width=0.85\textwidth]{slides-images/slide-19.jpg}
\caption{如何利用 Q-function 制定策略：选择使 $Q(s,a)$ 最大的动作\protect\footnotemark}
\end{figure}
\footnotetext{Slides 第19页。}

\subsection{本章小结}

强化学习的核心框架可以概括为一个闭环交互过程：Agent 观察当前状态 $s_t$，根据策略 $\pi$ 选择动作 $a_t$，环境返回新状态 $s_{t+1}$ 和奖励 $r_t$，Agent 的目标是找到使折扣累积回报最大化的最优策略 $\pi^*$。Q-function 是连接状态、动作和回报的桥梁，而策略是 Agent 的行为准则。


